\section{Related Work and Novelty Boundary}
\label{sec:related}

The relevant literature spans persistent agent memory, poisoning and backdoors, parameter-efficient personalization, provenance, unlearning, and federated robustness. The boundary matters: PoisonGuard is neither a new trigger-generation attack nor a federated aggregation rule. It is a lifecycle mechanism for a single device where one source item may influence two persistent sinks.

\subsection{Persistent memory and retrieval poisoning}

Retrieval-augmented generation combines parametric and retrieved knowledge \cite{lewis2020rag}. Agent systems extend this idea into longer-lived state: Generative Agents store and reflect over experience \cite{park2023generativeagents}, MemoryBank maintains long-term conversational memory \cite{zhong2023memorybank}, MemGPT manages context as an operating-system-like hierarchy \cite{packer2023memgpt}, and Reflexion converts feedback into reusable verbal memory \cite{shinn2023reflexion}. A survey by Zhang et al. organizes the growing design space of agent memory \cite{zhang2024memorysurvey}. These systems motivate persistent personalization but do not make mobile channel origin an enforceable training policy.

PoisonedRAG demonstrates that an attacker can corrupt retrieval by injecting carefully chosen texts \cite{zou2024poisonedrag}. AgentPoison optimizes triggers that retrieve malicious demonstrations from an agent memory or knowledge base \cite{chen2024agentpoison}. MINJA injects malicious records through query-only interaction \cite{dong2025minja}. Later work broadens the persistence surface. MemoryGraft targets stored experiences \cite{srivastava2025memorygraft}; Zombie Agents turns one-time external exposure into self-reinforcing memory \cite{yang2026zombie}; GhostWriter plants and later activates memory payloads in tool-using personal agents \cite{torres2026ghostwriter}; and MemGhost studies stealthy single-email delivery into persistent personal agents \cite{zhang2026memghost}. SPORE addresses extraction rather than integrity, but underscores that memory interfaces create cross-session attack paths \cite{gao2026spore}.

SkillJack supplies the closest empirical analogy to laundering: poisoned trajectories are transformed into reusable skills, become less detectable, and can survive deletion of the original records \cite{ying2026skilljack}. PoisonGuard transfers this question to mobile parameter updates and makes the source-to-weight dependency explicit. It does not claim that SkillJack already evaluated LoRA weights; it did not.

\subsection{Memory and prompt-injection defenses}

Recent defenses filter, isolate, or validate persistent content. A trust-score and temporal-decay defense studies realistic memory-bank contamination \cite{sunil2026memorypoison}. A-MemGuard combines consensus-based validation with a dual-memory design \cite{wei2025amemguard}. AgentSys isolates worker contexts and restricts what returns to a privileged context \cite{wen2026agentsys}. DASGuard traces control-like content in files back to its origin before committing state \cite{tan2026dasguard}. These works support the premise that origin and lifecycle control matter, while leaving open the specific relation between mobile ingestion, training-data selection, and adapter checkpoints.

Prompt-injection defenses address a neighboring boundary. Spotlighting marks untrusted data so a model can distinguish it from instructions \cite{hines2024spotlighting}; StruQ separates instructions and data through structured queries \cite{chen2024struq}; AgentDojo supplies realistic tool-use tasks and attack cases \cite{debenedetti2024agentdojo}. PoisonGuard can host these checks inside its audit gate, but its primary invariant is different: no derived persistent state may lose the provenance of an influencing event.

\subsection{Adapter poisoning and textual backdoors}

Modern candidate SLMs inherit the Transformer architecture \cite{vaswani2017attention}. LoRA makes adaptation compact by representing a weight update as a low-rank product \cite{hu2021lora}; QLoRA couples low-rank adaptation with quantization \cite{dettmers2023qlora}. LoRATK shows that malicious LoRAs can retain backdoors when merged with benign task adapters \cite{liu2024loratk}. Token-level adapter experiments show that poisoned local training data can create backdoors while leaving clean task performance largely unchanged, and that behavioral and weight-level statistics offer complementary signals \cite{lelle2026token}. Gradient Assembly Poisoning targets the factorization itself in distributed LoRA systems \cite{dong2026gap}; its federated attacker differs from PoisonGuard's content-origin attacker, so we use it only to motivate factor-aware auditing.

The broader backdoor literature establishes that small poisoned subsets can implant trigger-conditioned behavior. BadNets identifies supply-chain vulnerability \cite{gu2017badnets}; targeted and clean-label attacks show that labels need not be visibly inconsistent \cite{chen2017targetedbackdoor,shafahi2018poisonfrogs}; and weight poisoning targets pretrained NLP models \cite{kurita2020weightpoison}. NeuBA operates at neuron level in pretrained language models \cite{zhang2021neuba}. Detection methods include spectral signatures \cite{tran2018spectral}, activation clustering \cite{chen2018activation}, influence functions \cite{koh2017influence}, STRIP's input perturbations \cite{gao2019strip}, and textual defenses such as ONION and RAP \cite{qi2020onion,yang2021rap}. PoisonGuard does not assume any one detector is complete. It uses detection to decide promotion and uses lineage for recovery after detection.

\subsection{On-device personalization}

MobileBERT demonstrated that language models can be engineered for resource-constrained inference \cite{sun2020mobilebert}. PocketLLM evaluates derivative-free local fine-tuning on a smartphone \cite{peng2024pocketllm}; MobileFineTuner provides a mobile-native framework for full and parameter-efficient training \cite{geng2025mobilefinetuner}; XPerT reduces local personalization cost through model selection \cite{wang2025xpert}; and recent hypernetwork work synthesizes personalized LoRAs on-device \cite{augenstein2026hyperlora}. LaMP supplies personalization tasks \cite{salemi2023lamp}, while PersonaChat and AESLC provide dialogue and email-like text workloads \cite{zhang2018personachat,zhang2019aeslc}. Qwen2.5 and Gemma 3 define candidate compact model families \cite{qwen2024qwen,gemma2025gemma}. This literature establishes feasibility, not deployed prevalence. Accordingly, our evaluation claims a prototype study rather than universal production coverage.

\subsection{Unlearning, rollback, and provenance}

Machine unlearning seeks to remove the effect of designated examples. SISA partitions training to limit retraining scope \cite{bourtoule2021machineunlearning}; certified removal provides model-specific guarantees \cite{guo2020certified}; adversarial evaluation exposes weaknesses of inexact methods \cite{goel2022adversarialunlearning}; and approximate LLM methods include task-specific forgetting and negative preference optimization \cite{eldan2023harry,zhang2024npo}. Most importantly for this work, Pawelczyk et al. find that practical approximate methods can fail to eliminate poisoning effects \cite{pawelczyk2025unlearning}. PoisonGuard therefore does not call a gradient-based forgetting step ``rollback.'' It restores a pre-influence checkpoint and re-executes a versioned clean suffix.

Data provenance traditionally records where outputs came from and how they were derived. PoisonGuard adopts that causal viewpoint but narrows it to stateful mobile learning. The lineage graph is not a blockchain and does not require distributed consensus. Its purpose is complete mediation, reachability-based invalidation, and deterministic replay.

\subsection{Federated robustness as an inapplicable contrast}

Federated averaging coordinates many clients through a server \cite{mcmahan2017fedavg}. Krum, coordinate-wise robust statistics, Bulyan, and geometric-median aggregation resist some Byzantine update models \cite{blanchard2017krum,yin2018byzantine,mhamdi2018bulyan,pillutla2019rfa}. FLTrust bootstraps trust from a server dataset \cite{cao2020fltrust}; FLAME clusters and clips updates before adding noise \cite{nguyen2021flame}; and FoolsGold addresses sybil-like similarity \cite{fung2018sybil}. Model-poisoning and backdoor studies illustrate the limits of aggregation \cite{bhagoji2018federated,bagdasaryan2020backdoorfl}. More recent federated instruction-tuning work shows that low-concentration poisoned data spread across nominally benign clients can evade defenses designed for malicious clients \cite{zhao2026fedbit}. These methods assume multiple updates to compare. A single phone learning from its own mixed-trust streams has neither that population nor a server root dataset. PoisonGuard cites them to delimit the problem, not as fair head-to-head baselines.

\begin{table}[t]
\centering
\small
\caption{Novelty boundary across related literature.}
\label{tab:boundary}
\resizebox{\columnwidth}{!}{%
\begin{tabular}{lccccc}
\toprule
Family & Memory & Adapter & Provenance & Cross-sink & Exact replay \\
\midrule
Memory/RAG attacks \cite{chen2024agentpoison,dong2025minja,zou2024poisonedrag} & \checkmark & -- & limited & -- & -- \\
Persistent-agent attacks \cite{zhang2026memghost,ying2026skilljack} & \checkmark & skill & motivates & partial & -- \\
Memory defenses \cite{sunil2026memorypoison,wei2025amemguard} & \checkmark & -- & trust signal & -- & -- \\
Adapter defenses \cite{tran2018spectral,lelle2026token} & -- & \checkmark & -- & -- & -- \\
Approximate unlearning \cite{eldan2023harry,zhang2024npo} & -- & \checkmark & deletion set & limited & -- \\
\PG{} (this work) & \checkmark & \checkmark & \checkmark & \checkmark & \checkmark \\
\bottomrule
\end{tabular}%
}
\end{table}
